\documentclass{article}
\usepackage{amsmath}

\title{A Cigányok története}
\author{Rovenszky Robin}
\date{Legutóbb frissítve: 2025. Március 13.}

\begin{document}

\maketitle

\section{Bevezető}
A 9-11. század között elhagyják Indiát, Perzsián és a Bizánczi Birodalmon keresztül érkeztek Európába a 14-15. század környékén.

A cigány szó görögből ered, érinthetetlent jelent. A roma szó a saját nyelvükből ered és férfit jelent.

\subsection{Eredeti foglalkozásaik}
A kora középkorban alkalmi munkák; teknővájás, muzsikálás, lókereskedés.

\subsection{Csoportjaik}

\begin{enumerate}
    \item Magyarcigányok (romungrók) $\rightarrow$ cigányzene. Kb 600-900 ezer roma, főleg Szabolcs-Szatmár-Bereg megyében.
    \item Oláh cigányok: Sokáig vándorló életmódot fojtattak, fenntartják a régi nyelvüket, törvényeiket.
    \item Beások: Archaikus csoport, hagyományos cigány foglalkozások a kora középkorból. 
\end{enumerate}

\subsection{Értékek, kultúra, hagyományok}
Tánc, zene, idősek felé való tisztelet.

\subsection{Történelem}
Parrajmos: cigány holokauszt\\
Kommunizmusban erőszakosan megpróbálta raáerőszakolni a cigányokra a rendszert, hogy letelepedjenek, munkát vállaljanak stb., majd a rendszerváltás után, ahogy megszűntek a munkalehetőségek, nagyon nehéz volt nekik munkát találniuk.\\

Kirekesztés: oktatási szegregáció; külön cigányiskola, rosszabb felszerelés $\implies$ nem megfelelő taníttatás $\implies$ nehéz kitörni, legfőképpen kétkezi munkák, gettó, nincs jó víz/gáz/fűtés, előítéletek / sztereotípiák.


\end{document}